\documentclass[11pt,a4paper]{article}
\PassOptionsToPackage{obeyspaces,spaces}{url}

\usepackage[margin=2.3cm]{geometry}
\usepackage[T1]{fontenc}
\usepackage{lmodern}
\usepackage{microtype}
\usepackage{xcolor}
\usepackage{graphicx}
\usepackage{booktabs}
\usepackage{tabularx}
\usepackage{enumitem}
\usepackage{listings}
\usepackage{algorithm}
\usepackage{algpseudocode}
\usepackage[most]{tcolorbox}
\usepackage{amssymb}
\usepackage{hyperref}
\usepackage{xurl}

\definecolor{accent}{HTML}{1F4E79}
\definecolor{codebg}{HTML}{F4F6F8}
\definecolor{good}{HTML}{2E7D32}
\definecolor{warn}{HTML}{B26A00}
\definecolor{bad}{HTML}{C62828}

\hypersetup{colorlinks=true, linkcolor=accent, urlcolor=accent}
\setlist{itemsep=2pt, topsep=4pt}

\lstset{
  basicstyle=\ttfamily\small,
  backgroundcolor=\color{codebg},
  frame=single, rulecolor=\color{codebg},
  breaklines=true, columns=fullflexible,
  xleftmargin=4pt, xrightmargin=4pt,
  showstringspaces=false
}

\newtcolorbox{keybox}[1]{colback=accent!5, colframe=accent, fonttitle=\bfseries, title=#1, arc=2pt, boxrule=0.6pt}
\newtcolorbox{warnbox}[1]{colback=warn!6, colframe=warn, fonttitle=\bfseries, title=#1, arc=2pt, boxrule=0.6pt}

\title{\textbf{Pacman Capture the Flag}\\[4pt]\large Setup, Diagnosis and Three Competitive Strategies}
\author{Team \texttt{pacman-agent}}
\date{October 2026}

\begin{document}
\maketitle

\begin{abstract}
\noindent This document collects everything needed to compete in the Pacman Capture-the-Flag contest:
a digest of the rules that matter for strategy, the root cause of the macOS simulation problem and
how it was fixed, a cheat-sheet for running games locally, and three strategies of increasing strength
and risk. \texttt{my\_team.py} is still the organisers' original template. The strategies below describe what to build in it.
A prototype of Strategy~A won 40/40 test games against the baseline.
\end{abstract}

\tableofcontents
\newpage

%=====================================================================
\section{Rules That Shape Strategy}
%=====================================================================

The rules come from \texttt{CONTEST.md} and the constants in \texttt{capture.py}.

\begin{center}
\begin{tabularx}{\linewidth}{@{}lX@{}}
\toprule
\textbf{Rule} & \textbf{Consequence for our agents} \\
\midrule
Map split in two halves & You are a ghost (can eat) at home and a Pacman (can be eaten) on the enemy side. \\
Food scores only when brought home & Carrying food is risky. Bank it regularly instead of greedily eating everything. \\
Death drops carried food & Dying while carrying 10 dots throws away 10 points of work and gives the defender time. \\
Win condition & Return 28 dots ($\tfrac{60}{2}-2$), or lead after 1200 moves (300 per agent). \\
Ranking: win 3, tie 1, loss 0 & A 0--0 draw is nearly worthless. Pure defence is not enough. \\
Power capsule & Enemy ghosts are scared for 40 moves and can be eaten (they respawn at start). \\
Partial observability & Exact enemy positions only within Manhattan distance 5. Otherwise noisy distance $\pm 6$. \\
Time limits & 15\,s setup, 1\,s per move (3rd warning forfeits), any move $>3$\,s forfeits. \\
Hardware & 1.7\,GHz, 1.7\,GB RAM, so budget $\approx 0.3$\,s per move to be safe. \\
Layouts & Default plus random symmetric mazes. Either side may start, so never hard-code a map. \\
No threads, no printing & Threading is disqualifying. stdout/stderr is published with the results. \\
\bottomrule
\end{tabularx}
\end{center}

\begin{keybox}{Only one file is submitted}
The tournament runs \texttt{python3 capture.py -r teams/<a>/my\_team.py -b teams/<b>/my\_team.py}.
\texttt{create\_team()} is called \emph{without} extra options, so its default arguments are the team
that plays. Helper modules must be imported after appending the file's own folder to \texttt{sys.path}.
\end{keybox}

%=====================================================================
\section{Why the Simulation Did Not Work on the MacBook}
%=====================================================================

\subsection{Symptom}
Launching a game with graphics (\texttt{python capture.py ...}) never showed a usable window: the
process appeared frozen. Headless runs \emph{seemed} to hang too, because each game took around
40~seconds with almost no output.

\subsection{Diagnosis}
\begin{enumerate}
  \item The original \texttt{venv} was created from \texttt{/usr/bin/python3}. On macOS that is the
        \textbf{Python 3.9.6 bundled with Apple's Command Line Tools}
        (\texttt{venv/pyvenv.cfg}: \nolinkurl{home = /Library/Developer/CommandLineTools/usr/bin}).
  \item That Python links against Apple's \textbf{system Tcl/Tk 8.5} framework
        (\texttt{tkinter.TkVersion == 8.5}). Apple deprecated it years ago and no longer maintains it.
  \item The game's graphics layer (\texttt{graphics\_utils.py}) starts with \texttt{tkinter.Tk()}.
        A one-line probe, \texttt{tkinter.Tk().update()}, under that interpreter on
        \textbf{macOS 26.6} \emph{never returned}. It spun at 100\% CPU and grew to about
        1\,GB of RAM. The old Tk 8.5 cannot initialise its Cocoa window on current macOS.
  \item Homebrew's \textbf{Python 3.14 with Tk 9.1} (\texttt{python-tk@3.14}) was already installed.
        The same probe returned instantly, and a full graphical game ran correctly.
\end{enumerate}

\begin{keybox}{Root cause in one sentence}
The virtual environment was built on Apple's outdated Command-Line-Tools Python, whose Tk~8.5 GUI
toolkit hangs on macOS~26. The game code is fine; the interpreter's GUI library is broken.
\end{keybox}

\subsection{Secondary problems that made it look worse}
\begin{center}
\begin{tabularx}{\linewidth}{@{}lX@{}}
\toprule
\textbf{Problem} & \textbf{Explanation and fix} \\
\midrule
40\,s per ``headless'' game & \texttt{-{}-delay-step} defaults to 0.03\,s, and $1200 \times 0.03 = 36$\,s of pure sleep. Pass \texttt{-{}-delay-step 0}. A game then takes under 1\,s. \\
``No teams found!'' & \texttt{-r baseline\_team} is resolved relative to the \emph{current folder}, so it only works inside \texttt{src/contest}. Use full paths. \\
\texttt{scripts/play.sh} broken & It references \texttt{myTeam.py} (doesn't exist) and the broken \texttt{runner.py}. Use the README command instead. \\
\texttt{scripts/config.sh} & Re-adds the submodule from a \emph{different} fork. Do not run it. \\
\texttt{-{}-delay} in CONTEST.md & The real flag is \texttt{-{}-delay-step}. \\
\texttt{venv/lib 2} folder & A Finder/iCloud duplicate. The project lives on an iCloud-synced Desktop, which can corrupt virtual environments. Prefer e.g.\ \texttt{\textasciitilde/code/}. \\
\bottomrule
\end{tabularx}
\end{center}

\subsection{The fix (already applied)}
\begin{lstlisting}[language=bash]
cd pacman-agent
rm -rf venv
/opt/homebrew/bin/python3.14 -m venv venv     # Homebrew Python with Tk 9.1
source venv/bin/activate
pip install --upgrade pip
pip install -e pacman-contest                  # the 'contest' package
pip install frozendict numpy
python -c "import tkinter; print(tkinter.TkVersion)"   # prints 9.1
\end{lstlisting}
If Homebrew's Tk is ever missing: \texttt{brew install python@3.14 python-tk@3.14}.
Note that Homebrew \texttt{python@3.13} on this machine has \emph{no} Tk module.

%=====================================================================
\section{Running Simulations Locally}
%=====================================================================

As in the README: activate the venv, then run everything from the engine folder
(\nolinkurl{source venv/bin/activate}, then \nolinkurl{cd pacman-contest/src/contest}).

\begin{center}
\begin{tabularx}{\linewidth}{@{}lX@{}}
\toprule
\textbf{Goal} & \textbf{Command} \\
\midrule
Watch a game & \nolinkurl{python capture.py -r baseline_team -b ../../../my_team.py} \\
Random map & add \texttt{-l RANDOM13} \\
Other map & add \texttt{-l officeCapture} (any name in \texttt{layouts/}) \\
Slow motion & add \texttt{-{}-delay-step 0.1} \\
Fast, no graphics & add \texttt{-q -{}-delay-step 0} \\
Ten games in a row & add \texttt{-n 10 -q -{}-delay-step 0} \\
Your team as blue & \nolinkurl{python capture.py -r ../../../my_team.py -b baseline_team} \\
Record a game & add \texttt{-{}-record}, then \nolinkurl{--replay www/contest_default/replays/match_0.replay} \\
Play yourself & add \texttt{-{}-keys0} (arrow keys / WASD control red agent 0) \\
All options & \texttt{python capture.py -{}-help} \\
\bottomrule
\end{tabularx}
\end{center}

\begin{warnbox}{Testing discipline}
\begin{itemize}
  \item Always test as \textbf{both red and blue} and on \textbf{random layouts}. Use \texttt{-l RANDOM} and swap \texttt{-r}/\texttt{-b}.
  \item Before replacing \texttt{my\_team.py}, keep a copy of the current one and play the new version against it.
  \item Keep a few strong opponent versions around. Beating only the baseline proves little.
\end{itemize}
\end{warnbox}

%=====================================================================
\section{Strategy Overview}
%=====================================================================

\begin{center}
\begin{tabularx}{\linewidth}{@{}lXXX@{}}
\toprule
 & \textbf{A: Smart Reflex} & \textbf{B: Tracking + Search} & \textbf{C: Adversarial Search} \\
\midrule
Core idea & Rule-based goals + Dijkstra paths & A + belief tracking + A* intercepts & B + alpha-beta / MCTS in close fights \\
Strength & Solid & Strong & Highest ceiling \\
Effort & Low (done) & Medium & High \\
Timeout risk & \textcolor{good}{None (7\,ms)} & \textcolor{good}{Low} & \textcolor{warn}{Needs time guards} \\
Main weakness & Blind to unseen enemies & Hand-tuned thresholds & Hard to tune, slow \\
\bottomrule
\end{tabularx}
\end{center}

\noindent\textbf{Recommended path:} submit A now, then layer B on top, and add C only for
close-combat situations. Benchmark every step against the previous version.

%=====================================================================
\section{Strategy A --- ``Smart Reflex''}
%=====================================================================

\subsection{Idea}
One \textbf{attacker} and one \textbf{defender}. Each turn, each agent picks a \emph{goal} with simple
rules, then plans the cheapest path to it with Dijkstra. Cells near dangerous ghosts get a heavy
cost, and cells adjacent to a ghost are blocked. Everything that can be precomputed is computed once
in \texttt{register\_initial\_state}.

\subsection{Map analysis (precomputed)}
\begin{itemize}
  \item \textbf{Home border:} the column of our cells next to the middle line. Going home means reaching
        \emph{any} border cell, not the spawn point.
  \item \textbf{Dead ends:} cells with one open neighbour are pruned repeatedly. Each pruned cell
        gets a \emph{depth} (moves to escape the pocket) and an \emph{entrance} (the exit cell).
\end{itemize}

\subsection{Attacker}
\begin{algorithm}[H]
\caption{Attacker goal selection (each turn)}
\begin{algorithmic}[1]
\If{food left $\le 2$ \textbf{or} (carrying $>0$ \textbf{and} moves left $\le$ distance home $+2$)}
  \State \Return path to home border
\EndIf
\If{a dangerous ghost is within 5}
  \If{some capsule is closer to us than to every ghost} \State \Return path to that capsule \EndIf
  \If{carrying $>0$ \textbf{or} ghost within 2} \State \Return path to home border \EndIf
\EndIf
\If{carrying $\ge$ limit (6, or 4 when home is far)} \State \Return path to home border \EndIf
\State $S \gets$ food where (no ghost is closer than us) \textbf{and}
       (if in a pocket: we can reach it and escape before a ghost reaches the entrance)
\State Remove food near the teammate's target (if both attack)
\State \Return Dijkstra path to nearest food in $S$, with ghost-proximity costs
\end{algorithmic}
\end{algorithm}

\noindent A ghost counts as \emph{dangerous} unless its scared timer exceeds its distance to us.
While capsule power lasts, the attacker ignores those ghosts and eats greedily.

\subsection{Defender}
\begin{itemize}
  \item \textbf{Visible invader:} chase it along home cells only, so the defender never becomes a Pacman.
  \item \textbf{Scared:} shadow the invader at distance 2 instead of touching it, then strike when the timer ends.
  \item \textbf{Invisible invader:} food that vanished since the last observation reveals where it is. Go there.
  \item \textbf{No invader:} patrol the border cell with the smallest average distance to the food and capsules we still defend.
\end{itemize}

\subsection{Role switching}
\begin{itemize}
  \item The attacker, while still a ghost at home, eats any invader within 4 steps (``free kill'').
  \item Late game (fewer than 60 moves each) and ahead: the attacker drops back to defend the lead.
  \item Second half, not winning, and no invaders: the defender joins the attack.
\end{itemize}

\subsection{Prototype results}
\begin{center}
\begin{tabular}{@{}lcccc@{}}
\toprule
Opponent & Games & W / T / L & Avg.\ margin & Max move time \\
\midrule
\texttt{baseline\_team}$^{*}$ & 40 & \textcolor{good}{40 / 0 / 0} & +14 & 7\,ms \\
\bottomrule
\end{tabular}\\[2pt]
{\footnotesize $^{*}$10 fixed + 10 random layouts, each played as red and as blue.}
\end{center}

\subsection{Tuning knobs}
\texttt{CARRY\_LIMIT} (6), \texttt{DANGER\_RADIUS} (5), \texttt{SAFETY\_MARGIN} (2),\\ the ghost cost
weights in \texttt{ghost\_cost\_fn}, and the role-switch thresholds (60 / 200 moves). Change one at a time
and re-test against the previous version.

%=====================================================================
\section{Strategy B --- Tracking + Search}
%=====================================================================

\subsection{Idea}
Strategy A is blind beyond distance 5. Strategy B maintains a \textbf{probability distribution over
each enemy's position}, so both agents act on where enemies \emph{probably} are.

\subsection{Belief tracking (exact inference over cells)}
\begin{algorithm}[H]
\caption{Per-enemy belief update (each turn)}
\begin{algorithmic}[1]
\State \textbf{Init:} all mass on the enemy's start cell
\State \textbf{Elapse time:} spread each cell's mass uniformly over itself and its neighbours
\State \textbf{Observe:} multiply by $P(\text{noisy dist}\mid \text{true dist})$ via \texttt{get\_distance\_prob}
\State Zero out cells inconsistent with \texttt{is\_pacman} (which side the enemy is on)
\State Zero out cells within 5 of us or our teammate (we would see it there)
\If{enemy visible} \State put all mass on its position \EndIf
\If{our food vanished} \State put mass on that food cell \EndIf
\If{enemy just died (position $\to$ start)} \State reset to start \EndIf
\State Normalise (if all zero, reset to uniform over its side)
\end{algorithmic}
\end{algorithm}
\noindent Cost is $O(\text{cells})$ per enemy per turn, a few milliseconds. Both agents share one belief
through a module-level dictionary.

\subsection{Uses of the belief}
\begin{itemize}
  \item \textbf{Defender intercepts:} instead of chasing from behind, move to the border chokepoint on
        the path between the invader's most likely cell and our densest food.
  \item \textbf{Attacker safety:} treat high-probability cells as soft ghost positions in the Dijkstra cost,
        so it avoids ghosts it cannot see yet.
  \item \textbf{Entry point choice:} cross the border where expected ghost distance is largest.
  \item \textbf{Food clustering:} plan ``eat $k$ nearby dots, then escape'' routes with A*,
        scoring $\frac{\text{dots}}{\text{route length}}$ minus risk.
  \item \textbf{Two-agent split:} assign food clusters (e.g.\ top/bottom half) so attackers never compete.
\end{itemize}

\subsection{Pros and risks}
Large gain against mid-level teams that slip in unseen. The risk is mostly bugs (beliefs collapsing to
zero), so keep the reset rule and visualise beliefs with \texttt{self.debug\_draw}.

%=====================================================================
\section{Strategy C --- Adversarial Search / MCTS}
%=====================================================================

\subsection{Idea}
Reflex rules handle most of the game, but life-or-death moments (escaping a ghost, cornering a Pacman)
are where games are decided. When an enemy is within about 6 steps, switch to \textbf{look-ahead search}
over the real game model (\texttt{game\_state.generate\_successor}).

\subsection{Algorithm}
\begin{itemize}
  \item \textbf{Players:} us plus the nearest one or two enemies. The teammate follows its reflex policy.
  \item \textbf{Unseen enemies:} place them at their most likely belief cell (from B).
  \item \textbf{Search:} alpha-beta, iterative deepening (depth 2, 4, 6, ...), or MCTS with UCT for
        roughly 0.4\,s.
  \item \textbf{Leaf evaluation:} score, carried food, distance home, $-\infty$ if eaten while carrying,
        $+$ for eating an invader, food distance --- reuse A/B features.
  \item \textbf{Time guard:} check \texttt{time.time()} in the search loop and return the best move found
        so far. Fall back to the reflex move if no search completes.
\end{itemize}

\begin{algorithm}[H]
\caption{Hybrid controller}
\begin{algorithmic}[1]
\State $t_0 \gets$ now
\If{nearest relevant enemy (seen or believed) within 6}
  \State $a \gets$ iterative-deepening alpha-beta until now $- t_0 > 0.4$\,s
\Else
  \State $a \gets$ Strategy B action
\EndIf
\State \Return $a$
\end{algorithmic}
\end{algorithm}

\subsection{Pros and risks}
Highest skill ceiling. The searcher finds escapes and traps that hand-written rules miss. Risks: the
grading machine is slower than a MacBook (budget conservatively), \texttt{generate\_successor} is
expensive (cache successors), and the evaluation needs careful tuning. Test with the \texttt{-i} option
and watch the output for \texttt{took too long} warnings.

%=====================================================================
\section{Roadmap and Submission Checklist}
%=====================================================================

\subsection{Roadmap}
\begin{enumerate}
  \item \textbf{First:} build Strategy A in \texttt{my\_team.py}. It is safe and beats the baseline.
  \item Keep a copy of it. Every new version must beat it as red and as blue.
  \item Add belief tracking (B), visualise it, and use it in the defender first, where the gain is biggest.
  \item Add belief-aware attacker costs and food-cluster routes.
  \item Add time-boxed search (C) for close encounters only.
  \item Study tournament replays of losses. Each loss usually points at one missing rule.
\end{enumerate}

\subsection{Before every submission}
\begin{itemize}[label=$\square$]
  \item No \texttt{print} statements or debug drawing left in \texttt{my\_team.py}.
  \item \texttt{create\_team} defaults select the intended agents.
  \item Many games with \texttt{-n 10 -l RANDOM}: zero errors, zero time warnings.
  \item Works as both red and blue, on random maps.
  \item \texttt{TEAM.md} filled in.
  \item Any helper modules imported via the \texttt{sys.path} trick.
\end{itemize}

\vfill
\begin{center}
\includegraphics[width=0.55\linewidth]{../pacman-contest/src/contest/img/contestLayout.png}
\end{center}

\end{document}
